% Capítulo 3: archivo incluido desde main.tex mediante % Capítulo 3: archivo incluido desde main.tex mediante % Capítulo 3: archivo incluido desde main.tex mediante % Capítulo 3: archivo incluido desde main.tex mediante \input{chapters/03_estado_arte}.
\chapter{Estado del arte}
\label{chap:estado-arte}

Este capítulo revisa los conceptos y tecnologías que fundamentan el proyecto:
agentes conversacionales con recuperación de conocimiento, plataformas de
orquestación, integración web y gobierno de la IA generativa. El objetivo no es
atribuir al artefacto capacidades que no hayan sido verificadas, sino justificar
las decisiones de diseño y situar sus límites.

\section{De los asistentes basados en reglas a los agentes fundamentados}
\label{sec:evolucion-agentes}

Los asistentes conversacionales clásicos suelen organizarse como árboles de
diálogo o conjuntos de intenciones. Esta estrategia puede resultar predecible
para procesos muy acotados, pero exige anticipar numerosas formulaciones y rutas
de conversación. En una consulta de ciberseguridad, una misma intención puede
expresarse de forma muy distinta: por ejemplo, un usuario puede preguntar por un
correo sospechoso, un enlace que no reconoce o un posible robo de credenciales.
El problema se agrava cuando la respuesta depende de una política que cambia, de
un procedimiento interno o del contexto de un incidente.

Los modelos de lenguaje permiten generar una respuesta a partir de lenguaje
natural sin definir cada formulación de forma manual. Esta capacidad facilita el
acceso a documentación heterogénea, pero no convierte al modelo en una fuente de
verdad. Los modelos generativos pueden producir contenido erróneo presentado con
confianza, un fenómeno denominado por NIST \emph{confabulación}
\cite{nistGenAi2024}. Para un dominio como la ciberseguridad, la calidad no debe
medirse sólo por la fluidez de la conversación: importa que los hechos estén
fundamentados, que las recomendaciones sean prudentes y que la respuesta no
amplíe una solicitud peligrosa.

La generación aumentada por recuperación o RAG incorpora una fuente documental
externa al proceso de generación. Lewis et al. describen arquitecturas que
combinan memoria paramétrica y no paramétrica para tareas intensivas en
conocimiento, con el objetivo de recuperar pasajes relevantes y condicionar la
generación sobre ellos \cite{lewis2020rag}. En una solución empresarial, la idea
central es usar documentos y repositorios controlados como contexto de respuesta,
con la ventaja de poder actualizar el conocimiento sin reentrenar el modelo.

No obstante, RAG no garantiza por sí solo veracidad, actualidad o seguridad. La
recuperación puede seleccionar material irrelevante; el material puede estar
desactualizado, ser ambiguo o contener instrucciones maliciosas; y el modelo
puede resumirlo incorrectamente. Por ello, el presente trabajo trata el
conocimiento como un activo gobernado: se selecciona, se versiona, se limita por
alcance y se valida mediante preguntas cuya respuesta pueda contrastarse con una
fuente aprobada.

\section{Conocimiento de ciberseguridad para el caso de uso}
\label{sec:conocimiento-ciberseguridad}

El agente no necesita reproducir un marco completo de gestión de seguridad, pero
sí debe encajar en él. CSF 2.0 ofrece un vocabulario útil para relacionar las
consultas de personas usuarias con actividades de gobierno, protección,
detección, respuesta y recuperación \cite{nistCsf2024}. Por ejemplo, una duda
sobre autenticación multifactor se conecta principalmente con la protección; una
notificación de correo malicioso debe llevar hacia detección y respuesta; y una
pregunta sobre responsables o procedimiento remite a gobierno.

El contenido de respuesta ante incidentes requiere un tratamiento especialmente
cuidadoso. NIST recomienda incorporar la respuesta a incidentes en las
actividades de gestión del riesgo durante todo el ciclo de vida
\cite{nistIncident2025}. La guía de INCIBE sobre respuesta a incidentes subraya
la necesidad de responsables, procedimientos de notificación, registro de
evidencias y revisión del plan \cite{incibeCrisis}. De estos documentos se
extrae una decisión de diseño: el agente puede explicar pasos iniciales de
seguridad y el mecanismo de escalado, pero no debe simular una investigación,
ordenar acciones irreversibles ni sustituir a la cadena de decisión humana.

Los recursos públicos también ayudan a fijar un lenguaje accesible. INCIBE
describe el phishing como una suplantación destinada a obtener información,
causar un perjuicio económico o infectar el dispositivo \cite{incibePhishing}.
Esta definición sirve como punto de partida para preguntas frecuentes, pero la
respuesta del agente debe centrarse en la conducta segura: no abrir enlaces o
adjuntos sospechosos, usar el canal de reporte configurado y no aportar
credenciales ni datos personales en el chat.

La actualidad es otro requisito. El informe ENISA Threat Landscape 2025 analiza
incidentes y tendencias dentro de un periodo concreto y evidencia que la
prioridad de amenazas depende del contexto temporal y sectorial
\cite{enisaThreat2025}. Por tanto, no se deberá presentar una fuente de panorama
de amenazas como una descripción permanente ni inferir que sus estadísticas se
aplican a una organización ficticia. Su función en el TFM es justificar el
contexto y orientar la selección de casos, no alimentar afirmaciones específicas
que no estén respaldadas.

\section{Microsoft Copilot Studio como plataforma de agente}
\label{sec:copilot-studio}

Microsoft Copilot Studio es una plataforma para crear, configurar, probar y
publicar agentes con fuentes de conocimiento, temas, herramientas y canales. La
documentación de producto describe que el conocimiento puede añadirse a nivel de
agente o de tema y que admite, según el entorno y la licencia, sitios públicos,
documentos, SharePoint, Dataverse, Azure AI Search y otros orígenes
\cite{copilotKnowledge}. Esta diversidad es útil, pero aumenta la importancia de
aplicar el principio de mínimo conocimiento: se incorporará únicamente el
material necesario para las categorías de consulta aprobadas.

La orquestación generativa introduce una capa de planificación basada en modelos
de lenguaje que interpreta la intención y puede seleccionar conocimientos,
herramientas o temas para construir una respuesta \cite{copilotOrchestration}.
Frente a un árbol de diálogo totalmente manual, esa aproximación reduce parte de
la rigidez conversacional; a cambio, el comportamiento depende de las
instrucciones, descripciones y recursos disponibles. La plataforma recomienda
que las instrucciones se apoyen en recursos configurados, empleen nombres
exactos cuando se refieran a herramientas y se prueben tras cada cambio
\cite{copilotInstructions}.

Para este TFM, las instrucciones se utilizarán como una capa de orientación, no
como un control de seguridad único. Deben establecer el dominio permitido, el
formato de respuesta, la obligación de indicar límites y las condiciones de
derivación. Las decisiones críticas ---como bloquear la publicación no
autenticada, filtrar endpoints o restringir conectores--- requieren además
configuración de plataforma, control de acceso y revisión humana.

Las fuentes públicas merecen una precaución adicional. La documentación de
Microsoft indica que la recuperación sobre sitios públicos depende de contenido
indexado y posee límites de profundidad y de cobertura \cite{copilotPublicWeb}.
En consecuencia, una URL añadida como fuente no equivale a haber validado que
todas sus páginas relevantes serán recuperadas. El inventario de conocimiento
debe acompañarse de pruebas de recuperación para las preguntas que se pretenda
soportar.

\section{Evaluación de agentes conversacionales}
\label{sec:evaluacion-agentes}

La evaluación de un agente puede abarcar respuesta única, conversaciones de
varios turnos, recuperación de conocimiento, llamadas a herramientas, seguridad
y experiencia de usuario. Copilot Studio permite organizar casos en conjuntos de
prueba y comparar la salida con respuestas esperadas o métodos de calidad
\cite{copilotEvaluation}. Esta capacidad facilita la regresión: un mismo conjunto
puede ejecutarse después de modificar una fuente, una instrucción o una
integración.

Sin embargo, una evaluación automática debe interpretarse con cautela. Una
coincidencia textual puede penalizar una respuesta correcta formulada de otra
manera; una medida semántica puede aceptar una respuesta que omite una acción
esencial; y una puntuación de calidad no demuestra que la fuente recuperada sea
la adecuada. Por ello, el protocolo del Capítulo~\ref{chap:objetivos-metodologia}
combina criterios automáticos, cuando estén disponibles, y revisión humana con
una rúbrica. También reserva casos específicos para abstención segura,
derivación y resistencia frente a entradas adversariales.

La trazabilidad aporta un criterio adicional. Una respuesta no debe citar una
fuente inexistente, una URL no relacionada o una política ficticia como si fuera
real. Cuando el canal permita mostrar referencias, se comprobará que la
referencia corresponda al contenido utilizado. Cuando no las muestre, la
trazabilidad deberá aportarse mediante registros de prueba y un mapeo manual
entre respuesta esperada y fuente de conocimiento.

\section{Integración web y seguridad de canal}
\label{sec:integracion-web}

Una aplicación Angular permite controlar el contexto visual, los avisos al
usuario, los enlaces a procedimientos y la experiencia de consulta. Angular
recomienda mantener las dependencias actualizadas, evitar APIs identificadas como
riesgo de seguridad y reducir el riesgo de XSS mediante plantillas, sanitización,
Content Security Policy y Trusted Types cuando resulte aplicable
\cite{angularSecurity}. No obstante, la protección frente a XSS no resuelve por
sí sola los problemas propios de una aplicación de chat: gestión de sesión,
autorización, protección de tokens, controles de origen, registros y límites de
uso deben diseñarse de forma conjunta.

Copilot Studio puede publicar agentes en diversos canales. Microsoft recomienda
aplicar autenticación y controles de gobierno, especialmente cuando el agente se
utiliza dentro de una organización o por un grupo específico de usuarios
\cite{copilotPublish}. Para una aplicación web personalizada, el diseño deberá
declarar la opción de autenticación seleccionada y justificarla. Una integración
sin autenticación puede ser apropiada para una demostración sin conocimiento
interno, pero no constituye el diseño por defecto de un despliegue organizativo.

La documentación de seguridad del canal web indica que Copilot Studio utiliza el
canal Direct Line de Bot Framework para conectar la página o aplicación con el
agente. Si se habilita el acceso protegido, el cliente debe usar secretos o
tokens; en código que se ejecuta en el navegador, el secreto debe intercambiarse
por un token desde un servicio y no exponerse en el cliente
\cite{copilotWebSecurity}. Esta propiedad condiciona la arquitectura: un secreto
de Direct Line no se incluirá en variables públicas de Angular ni en un
repositorio.

Microsoft también documenta opciones de integración mediante SDK y Direct Line,
incluido un patrón basado en eventos para enviar contexto inicial de forma fiable
\cite{copilotDirectLine}. Dicho contexto deberá limitarse a datos no sensibles y
estar sujeto a validación. Si se habilita inicio de sesión único, Copilot Studio
requiere una configuración de autenticación con Microsoft Entra ID, registros de
aplicación diferenciados y un flujo de intercambio de tokens
\cite{copilotSso}. Esta configuración se tratará como una mejora opcional
condicionada a los requisitos y recursos del entorno, no como una funcionalidad
que pueda declararse implantada sin evidencia.

La viabilidad operativa también depende de licencias y capacidad. Microsoft
distingue licencias de usuario para crear y gestionar agentes de la licencia de
tenant, y recomienda estimar el consumo antes de un despliegue
\cite{copilotLicensing}. En este trabajo, los aspectos de licencia se documentan
como una restricción de entorno. No se infieren costes ni derechos de uso que no
puedan verificarse en el tenant concreto.

\section{Riesgos de seguridad específicos de la IA generativa}
\label{sec:riesgos-ia-generativa}

El modelo de amenazas debe considerar tanto vulnerabilidades tradicionales de la
aplicación como riesgos propios del agente. El AI RMF de NIST organiza la gestión
de riesgos de IA alrededor de las funciones de gobernar, mapear, medir y
gestionar \cite{nistAiRmf2023}. Su perfil para IA generativa señala riesgos
relacionados con la confabulación, privacidad, seguridad de la información,
integración de componentes y configuración humano--IA \cite{nistGenAi2024}.

OWASP incluye entre los riesgos de aplicaciones con LLM la inyección de prompts,
la divulgación de información sensible, los riesgos de cadena de suministro, el
envenenamiento de datos y la agencia excesiva \cite{owaspLlm2026}. Estos riesgos
se traducen en controles concretos para el proyecto:

\begin{itemize}
  \item no incorporar a la base de conocimiento texto no revisado ni permitir que
  una instrucción incluida en un documento redefina el comportamiento del agente;
  \item no almacenar secretos, tokens, credenciales o evidencias reales en
  fuentes de conocimiento ni en las conversaciones de prueba;
  \item limitar las acciones del agente a orientación informativa; no se
  habilitarán herramientas con capacidad de modificar sistemas;
  \item validar que las salidas presentadas por la interfaz no se interpreten como
  HTML confiable ni como instrucciones ejecutables;
  \item incorporar avisos visibles sobre el alcance, la posible incertidumbre y
  la necesidad de recurrir a canales humanos ante incidentes; y
  \item probar explícitamente intentos de revelar instrucciones, alterar el
  contexto o conseguir información de otros usuarios.
\end{itemize}

La defensa no debe descansar sobre una única instrucción de rechazo. La reducción
de superficie se consigue combinando alcance limitado, conocimiento aprobado,
autenticación adecuada, seguridad de canal, mínimo privilegio, revisión de
configuración y pruebas de robustez. Las limitaciones que sobrevivan a estos
controles deben aparecer en la memoria y en la interfaz, en lugar de ocultarse
detrás de un lenguaje excesivamente seguro.

\section{Privacidad, regulación y supervisión humana}
\label{sec:privacidad-regulacion}

Un agente conversacional puede recibir datos personales o información sensible
aunque no se haya diseñado para ello. El RGPD exige que el responsable aplique
medidas técnicas y organizativas apropiadas al riesgo y pueda demostrar el
cumplimiento \cite{gdpr}. Para el TFM, el control más directo es no utilizar
información real en pruebas y minimizar los datos que el usuario puede aportar.
En un uso futuro, habría que analizar la base jurídica, información al usuario,
retención, transferencias, subencargados, derechos de los interesados y gestión
de brechas.

El Reglamento (UE) 2024/1689 establece reglas armonizadas sobre sistemas de IA y
persigue promover una IA centrada en las personas y fiable, protegiendo a la vez
salud, seguridad y derechos fundamentales \cite{aiAct}. La clasificación y las
obligaciones concretas dependen del caso de uso, rol y fecha de aplicación; no se
deduce automáticamente que este prototipo cumpla o incumpla una categoría
regulatoria. Sus principios sí informan la solución: transparencia acerca de la
interacción con un agente, supervisión humana, documentación y tratamiento
proporcional del riesgo.

\section{Síntesis y contribución propuesta}
\label{sec:sintesis-estado-arte}

La Tabla~\ref{tab:comparacion-enfoques} resume el espacio de diseño. No pretende
medir el rendimiento de una implementación concreta; identifica qué propiedad
aporta cada enfoque y qué limitación debe compensarse.

\begin{table}[H]
\centering
\caption{Comparación conceptual de enfoques de asistencia conversacional.}
\label{tab:comparacion-enfoques}
\begin{tabular}{p{0.24\textwidth}p{0.31\textwidth}p{0.34\textwidth}}
\hline
\textbf{Enfoque} & \textbf{Ventaja principal} & \textbf{Limitación relevante}\\
\hline
FAQ o árbol de diálogo & Comportamiento predecible para rutas predefinidas. &
Escasa flexibilidad ante paráfrasis y mantenimiento costoso de rutas.\\
Modelo generativo sin conocimiento controlado & Interacción flexible y lenguaje
natural. & Riesgo elevado de información no verificable o desactualizada.\\
Agente generativo con recuperación y gobierno & Puede unir consultas variadas con
fuentes seleccionadas y trazabilidad. & Depende de la calidad de la recuperación,
de las instrucciones y de controles adicionales.\\
\hline
\end{tabular}
\end{table}

La contribución propuesta del TFM es un artefacto de alcance restringido que
combina una base de conocimiento curada, reglas explícitas de seguridad,
integración Angular y una batería de evaluación orientada a evidencia. Su valor
no reside en afirmar que un agente sustituye a especialistas, sino en mostrar
cómo puede diseñarse un primer nivel de apoyo que mantenga la trazabilidad,
reconozca sus límites y pueda evolucionar mediante pruebas reproducibles.

.
\chapter{Estado del arte}
\label{chap:estado-arte}

Este capítulo revisa los conceptos y tecnologías que fundamentan el proyecto:
agentes conversacionales con recuperación de conocimiento, plataformas de
orquestación, integración web y gobierno de la IA generativa. El objetivo no es
atribuir al artefacto capacidades que no hayan sido verificadas, sino justificar
las decisiones de diseño y situar sus límites.

\section{De los asistentes basados en reglas a los agentes fundamentados}
\label{sec:evolucion-agentes}

Los asistentes conversacionales clásicos suelen organizarse como árboles de
diálogo o conjuntos de intenciones. Esta estrategia puede resultar predecible
para procesos muy acotados, pero exige anticipar numerosas formulaciones y rutas
de conversación. En una consulta de ciberseguridad, una misma intención puede
expresarse de forma muy distinta: por ejemplo, un usuario puede preguntar por un
correo sospechoso, un enlace que no reconoce o un posible robo de credenciales.
El problema se agrava cuando la respuesta depende de una política que cambia, de
un procedimiento interno o del contexto de un incidente.

Los modelos de lenguaje permiten generar una respuesta a partir de lenguaje
natural sin definir cada formulación de forma manual. Esta capacidad facilita el
acceso a documentación heterogénea, pero no convierte al modelo en una fuente de
verdad. Los modelos generativos pueden producir contenido erróneo presentado con
confianza, un fenómeno denominado por NIST \emph{confabulación}
\cite{nistGenAi2024}. Para un dominio como la ciberseguridad, la calidad no debe
medirse sólo por la fluidez de la conversación: importa que los hechos estén
fundamentados, que las recomendaciones sean prudentes y que la respuesta no
amplíe una solicitud peligrosa.

La generación aumentada por recuperación o RAG incorpora una fuente documental
externa al proceso de generación. Lewis et al. describen arquitecturas que
combinan memoria paramétrica y no paramétrica para tareas intensivas en
conocimiento, con el objetivo de recuperar pasajes relevantes y condicionar la
generación sobre ellos \cite{lewis2020rag}. En una solución empresarial, la idea
central es usar documentos y repositorios controlados como contexto de respuesta,
con la ventaja de poder actualizar el conocimiento sin reentrenar el modelo.

No obstante, RAG no garantiza por sí solo veracidad, actualidad o seguridad. La
recuperación puede seleccionar material irrelevante; el material puede estar
desactualizado, ser ambiguo o contener instrucciones maliciosas; y el modelo
puede resumirlo incorrectamente. Por ello, el presente trabajo trata el
conocimiento como un activo gobernado: se selecciona, se versiona, se limita por
alcance y se valida mediante preguntas cuya respuesta pueda contrastarse con una
fuente aprobada.

\section{Conocimiento de ciberseguridad para el caso de uso}
\label{sec:conocimiento-ciberseguridad}

El agente no necesita reproducir un marco completo de gestión de seguridad, pero
sí debe encajar en él. CSF 2.0 ofrece un vocabulario útil para relacionar las
consultas de personas usuarias con actividades de gobierno, protección,
detección, respuesta y recuperación \cite{nistCsf2024}. Por ejemplo, una duda
sobre autenticación multifactor se conecta principalmente con la protección; una
notificación de correo malicioso debe llevar hacia detección y respuesta; y una
pregunta sobre responsables o procedimiento remite a gobierno.

El contenido de respuesta ante incidentes requiere un tratamiento especialmente
cuidadoso. NIST recomienda incorporar la respuesta a incidentes en las
actividades de gestión del riesgo durante todo el ciclo de vida
\cite{nistIncident2025}. La guía de INCIBE sobre respuesta a incidentes subraya
la necesidad de responsables, procedimientos de notificación, registro de
evidencias y revisión del plan \cite{incibeCrisis}. De estos documentos se
extrae una decisión de diseño: el agente puede explicar pasos iniciales de
seguridad y el mecanismo de escalado, pero no debe simular una investigación,
ordenar acciones irreversibles ni sustituir a la cadena de decisión humana.

Los recursos públicos también ayudan a fijar un lenguaje accesible. INCIBE
describe el phishing como una suplantación destinada a obtener información,
causar un perjuicio económico o infectar el dispositivo \cite{incibePhishing}.
Esta definición sirve como punto de partida para preguntas frecuentes, pero la
respuesta del agente debe centrarse en la conducta segura: no abrir enlaces o
adjuntos sospechosos, usar el canal de reporte configurado y no aportar
credenciales ni datos personales en el chat.

La actualidad es otro requisito. El informe ENISA Threat Landscape 2025 analiza
incidentes y tendencias dentro de un periodo concreto y evidencia que la
prioridad de amenazas depende del contexto temporal y sectorial
\cite{enisaThreat2025}. Por tanto, no se deberá presentar una fuente de panorama
de amenazas como una descripción permanente ni inferir que sus estadísticas se
aplican a una organización ficticia. Su función en el TFM es justificar el
contexto y orientar la selección de casos, no alimentar afirmaciones específicas
que no estén respaldadas.

\section{Microsoft Copilot Studio como plataforma de agente}
\label{sec:copilot-studio}

Microsoft Copilot Studio es una plataforma para crear, configurar, probar y
publicar agentes con fuentes de conocimiento, temas, herramientas y canales. La
documentación de producto describe que el conocimiento puede añadirse a nivel de
agente o de tema y que admite, según el entorno y la licencia, sitios públicos,
documentos, SharePoint, Dataverse, Azure AI Search y otros orígenes
\cite{copilotKnowledge}. Esta diversidad es útil, pero aumenta la importancia de
aplicar el principio de mínimo conocimiento: se incorporará únicamente el
material necesario para las categorías de consulta aprobadas.

La orquestación generativa introduce una capa de planificación basada en modelos
de lenguaje que interpreta la intención y puede seleccionar conocimientos,
herramientas o temas para construir una respuesta \cite{copilotOrchestration}.
Frente a un árbol de diálogo totalmente manual, esa aproximación reduce parte de
la rigidez conversacional; a cambio, el comportamiento depende de las
instrucciones, descripciones y recursos disponibles. La plataforma recomienda
que las instrucciones se apoyen en recursos configurados, empleen nombres
exactos cuando se refieran a herramientas y se prueben tras cada cambio
\cite{copilotInstructions}.

Para este TFM, las instrucciones se utilizarán como una capa de orientación, no
como un control de seguridad único. Deben establecer el dominio permitido, el
formato de respuesta, la obligación de indicar límites y las condiciones de
derivación. Las decisiones críticas ---como bloquear la publicación no
autenticada, filtrar endpoints o restringir conectores--- requieren además
configuración de plataforma, control de acceso y revisión humana.

Las fuentes públicas merecen una precaución adicional. La documentación de
Microsoft indica que la recuperación sobre sitios públicos depende de contenido
indexado y posee límites de profundidad y de cobertura \cite{copilotPublicWeb}.
En consecuencia, una URL añadida como fuente no equivale a haber validado que
todas sus páginas relevantes serán recuperadas. El inventario de conocimiento
debe acompañarse de pruebas de recuperación para las preguntas que se pretenda
soportar.

\section{Evaluación de agentes conversacionales}
\label{sec:evaluacion-agentes}

La evaluación de un agente puede abarcar respuesta única, conversaciones de
varios turnos, recuperación de conocimiento, llamadas a herramientas, seguridad
y experiencia de usuario. Copilot Studio permite organizar casos en conjuntos de
prueba y comparar la salida con respuestas esperadas o métodos de calidad
\cite{copilotEvaluation}. Esta capacidad facilita la regresión: un mismo conjunto
puede ejecutarse después de modificar una fuente, una instrucción o una
integración.

Sin embargo, una evaluación automática debe interpretarse con cautela. Una
coincidencia textual puede penalizar una respuesta correcta formulada de otra
manera; una medida semántica puede aceptar una respuesta que omite una acción
esencial; y una puntuación de calidad no demuestra que la fuente recuperada sea
la adecuada. Por ello, el protocolo del Capítulo~\ref{chap:objetivos-metodologia}
combina criterios automáticos, cuando estén disponibles, y revisión humana con
una rúbrica. También reserva casos específicos para abstención segura,
derivación y resistencia frente a entradas adversariales.

La trazabilidad aporta un criterio adicional. Una respuesta no debe citar una
fuente inexistente, una URL no relacionada o una política ficticia como si fuera
real. Cuando el canal permita mostrar referencias, se comprobará que la
referencia corresponda al contenido utilizado. Cuando no las muestre, la
trazabilidad deberá aportarse mediante registros de prueba y un mapeo manual
entre respuesta esperada y fuente de conocimiento.

\section{Integración web y seguridad de canal}
\label{sec:integracion-web}

Una aplicación Angular permite controlar el contexto visual, los avisos al
usuario, los enlaces a procedimientos y la experiencia de consulta. Angular
recomienda mantener las dependencias actualizadas, evitar APIs identificadas como
riesgo de seguridad y reducir el riesgo de XSS mediante plantillas, sanitización,
Content Security Policy y Trusted Types cuando resulte aplicable
\cite{angularSecurity}. No obstante, la protección frente a XSS no resuelve por
sí sola los problemas propios de una aplicación de chat: gestión de sesión,
autorización, protección de tokens, controles de origen, registros y límites de
uso deben diseñarse de forma conjunta.

Copilot Studio puede publicar agentes en diversos canales. Microsoft recomienda
aplicar autenticación y controles de gobierno, especialmente cuando el agente se
utiliza dentro de una organización o por un grupo específico de usuarios
\cite{copilotPublish}. Para una aplicación web personalizada, el diseño deberá
declarar la opción de autenticación seleccionada y justificarla. Una integración
sin autenticación puede ser apropiada para una demostración sin conocimiento
interno, pero no constituye el diseño por defecto de un despliegue organizativo.

La documentación de seguridad del canal web indica que Copilot Studio utiliza el
canal Direct Line de Bot Framework para conectar la página o aplicación con el
agente. Si se habilita el acceso protegido, el cliente debe usar secretos o
tokens; en código que se ejecuta en el navegador, el secreto debe intercambiarse
por un token desde un servicio y no exponerse en el cliente
\cite{copilotWebSecurity}. Esta propiedad condiciona la arquitectura: un secreto
de Direct Line no se incluirá en variables públicas de Angular ni en un
repositorio.

Microsoft también documenta opciones de integración mediante SDK y Direct Line,
incluido un patrón basado en eventos para enviar contexto inicial de forma fiable
\cite{copilotDirectLine}. Dicho contexto deberá limitarse a datos no sensibles y
estar sujeto a validación. Si se habilita inicio de sesión único, Copilot Studio
requiere una configuración de autenticación con Microsoft Entra ID, registros de
aplicación diferenciados y un flujo de intercambio de tokens
\cite{copilotSso}. Esta configuración se tratará como una mejora opcional
condicionada a los requisitos y recursos del entorno, no como una funcionalidad
que pueda declararse implantada sin evidencia.

La viabilidad operativa también depende de licencias y capacidad. Microsoft
distingue licencias de usuario para crear y gestionar agentes de la licencia de
tenant, y recomienda estimar el consumo antes de un despliegue
\cite{copilotLicensing}. En este trabajo, los aspectos de licencia se documentan
como una restricción de entorno. No se infieren costes ni derechos de uso que no
puedan verificarse en el tenant concreto.

\section{Riesgos de seguridad específicos de la IA generativa}
\label{sec:riesgos-ia-generativa}

El modelo de amenazas debe considerar tanto vulnerabilidades tradicionales de la
aplicación como riesgos propios del agente. El AI RMF de NIST organiza la gestión
de riesgos de IA alrededor de las funciones de gobernar, mapear, medir y
gestionar \cite{nistAiRmf2023}. Su perfil para IA generativa señala riesgos
relacionados con la confabulación, privacidad, seguridad de la información,
integración de componentes y configuración humano--IA \cite{nistGenAi2024}.

OWASP incluye entre los riesgos de aplicaciones con LLM la inyección de prompts,
la divulgación de información sensible, los riesgos de cadena de suministro, el
envenenamiento de datos y la agencia excesiva \cite{owaspLlm2026}. Estos riesgos
se traducen en controles concretos para el proyecto:

\begin{itemize}
  \item no incorporar a la base de conocimiento texto no revisado ni permitir que
  una instrucción incluida en un documento redefina el comportamiento del agente;
  \item no almacenar secretos, tokens, credenciales o evidencias reales en
  fuentes de conocimiento ni en las conversaciones de prueba;
  \item limitar las acciones del agente a orientación informativa; no se
  habilitarán herramientas con capacidad de modificar sistemas;
  \item validar que las salidas presentadas por la interfaz no se interpreten como
  HTML confiable ni como instrucciones ejecutables;
  \item incorporar avisos visibles sobre el alcance, la posible incertidumbre y
  la necesidad de recurrir a canales humanos ante incidentes; y
  \item probar explícitamente intentos de revelar instrucciones, alterar el
  contexto o conseguir información de otros usuarios.
\end{itemize}

La defensa no debe descansar sobre una única instrucción de rechazo. La reducción
de superficie se consigue combinando alcance limitado, conocimiento aprobado,
autenticación adecuada, seguridad de canal, mínimo privilegio, revisión de
configuración y pruebas de robustez. Las limitaciones que sobrevivan a estos
controles deben aparecer en la memoria y en la interfaz, en lugar de ocultarse
detrás de un lenguaje excesivamente seguro.

\section{Privacidad, regulación y supervisión humana}
\label{sec:privacidad-regulacion}

Un agente conversacional puede recibir datos personales o información sensible
aunque no se haya diseñado para ello. El RGPD exige que el responsable aplique
medidas técnicas y organizativas apropiadas al riesgo y pueda demostrar el
cumplimiento \cite{gdpr}. Para el TFM, el control más directo es no utilizar
información real en pruebas y minimizar los datos que el usuario puede aportar.
En un uso futuro, habría que analizar la base jurídica, información al usuario,
retención, transferencias, subencargados, derechos de los interesados y gestión
de brechas.

El Reglamento (UE) 2024/1689 establece reglas armonizadas sobre sistemas de IA y
persigue promover una IA centrada en las personas y fiable, protegiendo a la vez
salud, seguridad y derechos fundamentales \cite{aiAct}. La clasificación y las
obligaciones concretas dependen del caso de uso, rol y fecha de aplicación; no se
deduce automáticamente que este prototipo cumpla o incumpla una categoría
regulatoria. Sus principios sí informan la solución: transparencia acerca de la
interacción con un agente, supervisión humana, documentación y tratamiento
proporcional del riesgo.

\section{Síntesis y contribución propuesta}
\label{sec:sintesis-estado-arte}

La Tabla~\ref{tab:comparacion-enfoques} resume el espacio de diseño. No pretende
medir el rendimiento de una implementación concreta; identifica qué propiedad
aporta cada enfoque y qué limitación debe compensarse.

\begin{table}[H]
\centering
\caption{Comparación conceptual de enfoques de asistencia conversacional.}
\label{tab:comparacion-enfoques}
\begin{tabular}{p{0.24\textwidth}p{0.31\textwidth}p{0.34\textwidth}}
\hline
\textbf{Enfoque} & \textbf{Ventaja principal} & \textbf{Limitación relevante}\\
\hline
FAQ o árbol de diálogo & Comportamiento predecible para rutas predefinidas. &
Escasa flexibilidad ante paráfrasis y mantenimiento costoso de rutas.\\
Modelo generativo sin conocimiento controlado & Interacción flexible y lenguaje
natural. & Riesgo elevado de información no verificable o desactualizada.\\
Agente generativo con recuperación y gobierno & Puede unir consultas variadas con
fuentes seleccionadas y trazabilidad. & Depende de la calidad de la recuperación,
de las instrucciones y de controles adicionales.\\
\hline
\end{tabular}
\end{table}

La contribución propuesta del TFM es un artefacto de alcance restringido que
combina una base de conocimiento curada, reglas explícitas de seguridad,
integración Angular y una batería de evaluación orientada a evidencia. Su valor
no reside en afirmar que un agente sustituye a especialistas, sino en mostrar
cómo puede diseñarse un primer nivel de apoyo que mantenga la trazabilidad,
reconozca sus límites y pueda evolucionar mediante pruebas reproducibles.

.
\chapter{Estado del arte}
\label{chap:estado-arte}

Este capítulo revisa los conceptos y tecnologías que fundamentan el proyecto:
agentes conversacionales con recuperación de conocimiento, plataformas de
orquestación, integración web y gobierno de la IA generativa. El objetivo no es
atribuir al artefacto capacidades que no hayan sido verificadas, sino justificar
las decisiones de diseño y situar sus límites.

\section{De los asistentes basados en reglas a los agentes fundamentados}
\label{sec:evolucion-agentes}

Los asistentes conversacionales clásicos suelen organizarse como árboles de
diálogo o conjuntos de intenciones. Esta estrategia puede resultar predecible
para procesos muy acotados, pero exige anticipar numerosas formulaciones y rutas
de conversación. En una consulta de ciberseguridad, una misma intención puede
expresarse de forma muy distinta: por ejemplo, un usuario puede preguntar por un
correo sospechoso, un enlace que no reconoce o un posible robo de credenciales.
El problema se agrava cuando la respuesta depende de una política que cambia, de
un procedimiento interno o del contexto de un incidente.

Los modelos de lenguaje permiten generar una respuesta a partir de lenguaje
natural sin definir cada formulación de forma manual. Esta capacidad facilita el
acceso a documentación heterogénea, pero no convierte al modelo en una fuente de
verdad. Los modelos generativos pueden producir contenido erróneo presentado con
confianza, un fenómeno denominado por NIST \emph{confabulación}
\cite{nistGenAi2024}. Para un dominio como la ciberseguridad, la calidad no debe
medirse sólo por la fluidez de la conversación: importa que los hechos estén
fundamentados, que las recomendaciones sean prudentes y que la respuesta no
amplíe una solicitud peligrosa.

La generación aumentada por recuperación o RAG incorpora una fuente documental
externa al proceso de generación. Lewis et al. describen arquitecturas que
combinan memoria paramétrica y no paramétrica para tareas intensivas en
conocimiento, con el objetivo de recuperar pasajes relevantes y condicionar la
generación sobre ellos \cite{lewis2020rag}. En una solución empresarial, la idea
central es usar documentos y repositorios controlados como contexto de respuesta,
con la ventaja de poder actualizar el conocimiento sin reentrenar el modelo.

No obstante, RAG no garantiza por sí solo veracidad, actualidad o seguridad. La
recuperación puede seleccionar material irrelevante; el material puede estar
desactualizado, ser ambiguo o contener instrucciones maliciosas; y el modelo
puede resumirlo incorrectamente. Por ello, el presente trabajo trata el
conocimiento como un activo gobernado: se selecciona, se versiona, se limita por
alcance y se valida mediante preguntas cuya respuesta pueda contrastarse con una
fuente aprobada.

\section{Conocimiento de ciberseguridad para el caso de uso}
\label{sec:conocimiento-ciberseguridad}

El agente no necesita reproducir un marco completo de gestión de seguridad, pero
sí debe encajar en él. CSF 2.0 ofrece un vocabulario útil para relacionar las
consultas de personas usuarias con actividades de gobierno, protección,
detección, respuesta y recuperación \cite{nistCsf2024}. Por ejemplo, una duda
sobre autenticación multifactor se conecta principalmente con la protección; una
notificación de correo malicioso debe llevar hacia detección y respuesta; y una
pregunta sobre responsables o procedimiento remite a gobierno.

El contenido de respuesta ante incidentes requiere un tratamiento especialmente
cuidadoso. NIST recomienda incorporar la respuesta a incidentes en las
actividades de gestión del riesgo durante todo el ciclo de vida
\cite{nistIncident2025}. La guía de INCIBE sobre respuesta a incidentes subraya
la necesidad de responsables, procedimientos de notificación, registro de
evidencias y revisión del plan \cite{incibeCrisis}. De estos documentos se
extrae una decisión de diseño: el agente puede explicar pasos iniciales de
seguridad y el mecanismo de escalado, pero no debe simular una investigación,
ordenar acciones irreversibles ni sustituir a la cadena de decisión humana.

Los recursos públicos también ayudan a fijar un lenguaje accesible. INCIBE
describe el phishing como una suplantación destinada a obtener información,
causar un perjuicio económico o infectar el dispositivo \cite{incibePhishing}.
Esta definición sirve como punto de partida para preguntas frecuentes, pero la
respuesta del agente debe centrarse en la conducta segura: no abrir enlaces o
adjuntos sospechosos, usar el canal de reporte configurado y no aportar
credenciales ni datos personales en el chat.

La actualidad es otro requisito. El informe ENISA Threat Landscape 2025 analiza
incidentes y tendencias dentro de un periodo concreto y evidencia que la
prioridad de amenazas depende del contexto temporal y sectorial
\cite{enisaThreat2025}. Por tanto, no se deberá presentar una fuente de panorama
de amenazas como una descripción permanente ni inferir que sus estadísticas se
aplican a una organización ficticia. Su función en el TFM es justificar el
contexto y orientar la selección de casos, no alimentar afirmaciones específicas
que no estén respaldadas.

\section{Microsoft Copilot Studio como plataforma de agente}
\label{sec:copilot-studio}

Microsoft Copilot Studio es una plataforma para crear, configurar, probar y
publicar agentes con fuentes de conocimiento, temas, herramientas y canales. La
documentación de producto describe que el conocimiento puede añadirse a nivel de
agente o de tema y que admite, según el entorno y la licencia, sitios públicos,
documentos, SharePoint, Dataverse, Azure AI Search y otros orígenes
\cite{copilotKnowledge}. Esta diversidad es útil, pero aumenta la importancia de
aplicar el principio de mínimo conocimiento: se incorporará únicamente el
material necesario para las categorías de consulta aprobadas.

La orquestación generativa introduce una capa de planificación basada en modelos
de lenguaje que interpreta la intención y puede seleccionar conocimientos,
herramientas o temas para construir una respuesta \cite{copilotOrchestration}.
Frente a un árbol de diálogo totalmente manual, esa aproximación reduce parte de
la rigidez conversacional; a cambio, el comportamiento depende de las
instrucciones, descripciones y recursos disponibles. La plataforma recomienda
que las instrucciones se apoyen en recursos configurados, empleen nombres
exactos cuando se refieran a herramientas y se prueben tras cada cambio
\cite{copilotInstructions}.

Para este TFM, las instrucciones se utilizarán como una capa de orientación, no
como un control de seguridad único. Deben establecer el dominio permitido, el
formato de respuesta, la obligación de indicar límites y las condiciones de
derivación. Las decisiones críticas ---como bloquear la publicación no
autenticada, filtrar endpoints o restringir conectores--- requieren además
configuración de plataforma, control de acceso y revisión humana.

Las fuentes públicas merecen una precaución adicional. La documentación de
Microsoft indica que la recuperación sobre sitios públicos depende de contenido
indexado y posee límites de profundidad y de cobertura \cite{copilotPublicWeb}.
En consecuencia, una URL añadida como fuente no equivale a haber validado que
todas sus páginas relevantes serán recuperadas. El inventario de conocimiento
debe acompañarse de pruebas de recuperación para las preguntas que se pretenda
soportar.

\section{Evaluación de agentes conversacionales}
\label{sec:evaluacion-agentes}

La evaluación de un agente puede abarcar respuesta única, conversaciones de
varios turnos, recuperación de conocimiento, llamadas a herramientas, seguridad
y experiencia de usuario. Copilot Studio permite organizar casos en conjuntos de
prueba y comparar la salida con respuestas esperadas o métodos de calidad
\cite{copilotEvaluation}. Esta capacidad facilita la regresión: un mismo conjunto
puede ejecutarse después de modificar una fuente, una instrucción o una
integración.

Sin embargo, una evaluación automática debe interpretarse con cautela. Una
coincidencia textual puede penalizar una respuesta correcta formulada de otra
manera; una medida semántica puede aceptar una respuesta que omite una acción
esencial; y una puntuación de calidad no demuestra que la fuente recuperada sea
la adecuada. Por ello, el protocolo del Capítulo~\ref{chap:objetivos-metodologia}
combina criterios automáticos, cuando estén disponibles, y revisión humana con
una rúbrica. También reserva casos específicos para abstención segura,
derivación y resistencia frente a entradas adversariales.

La trazabilidad aporta un criterio adicional. Una respuesta no debe citar una
fuente inexistente, una URL no relacionada o una política ficticia como si fuera
real. Cuando el canal permita mostrar referencias, se comprobará que la
referencia corresponda al contenido utilizado. Cuando no las muestre, la
trazabilidad deberá aportarse mediante registros de prueba y un mapeo manual
entre respuesta esperada y fuente de conocimiento.

\section{Integración web y seguridad de canal}
\label{sec:integracion-web}

Una aplicación Angular permite controlar el contexto visual, los avisos al
usuario, los enlaces a procedimientos y la experiencia de consulta. Angular
recomienda mantener las dependencias actualizadas, evitar APIs identificadas como
riesgo de seguridad y reducir el riesgo de XSS mediante plantillas, sanitización,
Content Security Policy y Trusted Types cuando resulte aplicable
\cite{angularSecurity}. No obstante, la protección frente a XSS no resuelve por
sí sola los problemas propios de una aplicación de chat: gestión de sesión,
autorización, protección de tokens, controles de origen, registros y límites de
uso deben diseñarse de forma conjunta.

Copilot Studio puede publicar agentes en diversos canales. Microsoft recomienda
aplicar autenticación y controles de gobierno, especialmente cuando el agente se
utiliza dentro de una organización o por un grupo específico de usuarios
\cite{copilotPublish}. Para una aplicación web personalizada, el diseño deberá
declarar la opción de autenticación seleccionada y justificarla. Una integración
sin autenticación puede ser apropiada para una demostración sin conocimiento
interno, pero no constituye el diseño por defecto de un despliegue organizativo.

La documentación de seguridad del canal web indica que Copilot Studio utiliza el
canal Direct Line de Bot Framework para conectar la página o aplicación con el
agente. Si se habilita el acceso protegido, el cliente debe usar secretos o
tokens; en código que se ejecuta en el navegador, el secreto debe intercambiarse
por un token desde un servicio y no exponerse en el cliente
\cite{copilotWebSecurity}. Esta propiedad condiciona la arquitectura: un secreto
de Direct Line no se incluirá en variables públicas de Angular ni en un
repositorio.

Microsoft también documenta opciones de integración mediante SDK y Direct Line,
incluido un patrón basado en eventos para enviar contexto inicial de forma fiable
\cite{copilotDirectLine}. Dicho contexto deberá limitarse a datos no sensibles y
estar sujeto a validación. Si se habilita inicio de sesión único, Copilot Studio
requiere una configuración de autenticación con Microsoft Entra ID, registros de
aplicación diferenciados y un flujo de intercambio de tokens
\cite{copilotSso}. Esta configuración se tratará como una mejora opcional
condicionada a los requisitos y recursos del entorno, no como una funcionalidad
que pueda declararse implantada sin evidencia.

La viabilidad operativa también depende de licencias y capacidad. Microsoft
distingue licencias de usuario para crear y gestionar agentes de la licencia de
tenant, y recomienda estimar el consumo antes de un despliegue
\cite{copilotLicensing}. En este trabajo, los aspectos de licencia se documentan
como una restricción de entorno. No se infieren costes ni derechos de uso que no
puedan verificarse en el tenant concreto.

\section{Riesgos de seguridad específicos de la IA generativa}
\label{sec:riesgos-ia-generativa}

El modelo de amenazas debe considerar tanto vulnerabilidades tradicionales de la
aplicación como riesgos propios del agente. El AI RMF de NIST organiza la gestión
de riesgos de IA alrededor de las funciones de gobernar, mapear, medir y
gestionar \cite{nistAiRmf2023}. Su perfil para IA generativa señala riesgos
relacionados con la confabulación, privacidad, seguridad de la información,
integración de componentes y configuración humano--IA \cite{nistGenAi2024}.

OWASP incluye entre los riesgos de aplicaciones con LLM la inyección de prompts,
la divulgación de información sensible, los riesgos de cadena de suministro, el
envenenamiento de datos y la agencia excesiva \cite{owaspLlm2026}. Estos riesgos
se traducen en controles concretos para el proyecto:

\begin{itemize}
  \item no incorporar a la base de conocimiento texto no revisado ni permitir que
  una instrucción incluida en un documento redefina el comportamiento del agente;
  \item no almacenar secretos, tokens, credenciales o evidencias reales en
  fuentes de conocimiento ni en las conversaciones de prueba;
  \item limitar las acciones del agente a orientación informativa; no se
  habilitarán herramientas con capacidad de modificar sistemas;
  \item validar que las salidas presentadas por la interfaz no se interpreten como
  HTML confiable ni como instrucciones ejecutables;
  \item incorporar avisos visibles sobre el alcance, la posible incertidumbre y
  la necesidad de recurrir a canales humanos ante incidentes; y
  \item probar explícitamente intentos de revelar instrucciones, alterar el
  contexto o conseguir información de otros usuarios.
\end{itemize}

La defensa no debe descansar sobre una única instrucción de rechazo. La reducción
de superficie se consigue combinando alcance limitado, conocimiento aprobado,
autenticación adecuada, seguridad de canal, mínimo privilegio, revisión de
configuración y pruebas de robustez. Las limitaciones que sobrevivan a estos
controles deben aparecer en la memoria y en la interfaz, en lugar de ocultarse
detrás de un lenguaje excesivamente seguro.

\section{Privacidad, regulación y supervisión humana}
\label{sec:privacidad-regulacion}

Un agente conversacional puede recibir datos personales o información sensible
aunque no se haya diseñado para ello. El RGPD exige que el responsable aplique
medidas técnicas y organizativas apropiadas al riesgo y pueda demostrar el
cumplimiento \cite{gdpr}. Para el TFM, el control más directo es no utilizar
información real en pruebas y minimizar los datos que el usuario puede aportar.
En un uso futuro, habría que analizar la base jurídica, información al usuario,
retención, transferencias, subencargados, derechos de los interesados y gestión
de brechas.

El Reglamento (UE) 2024/1689 establece reglas armonizadas sobre sistemas de IA y
persigue promover una IA centrada en las personas y fiable, protegiendo a la vez
salud, seguridad y derechos fundamentales \cite{aiAct}. La clasificación y las
obligaciones concretas dependen del caso de uso, rol y fecha de aplicación; no se
deduce automáticamente que este prototipo cumpla o incumpla una categoría
regulatoria. Sus principios sí informan la solución: transparencia acerca de la
interacción con un agente, supervisión humana, documentación y tratamiento
proporcional del riesgo.

\section{Síntesis y contribución propuesta}
\label{sec:sintesis-estado-arte}

La Tabla~\ref{tab:comparacion-enfoques} resume el espacio de diseño. No pretende
medir el rendimiento de una implementación concreta; identifica qué propiedad
aporta cada enfoque y qué limitación debe compensarse.

\begin{table}[H]
\centering
\caption{Comparación conceptual de enfoques de asistencia conversacional.}
\label{tab:comparacion-enfoques}
\begin{tabular}{p{0.24\textwidth}p{0.31\textwidth}p{0.34\textwidth}}
\hline
\textbf{Enfoque} & \textbf{Ventaja principal} & \textbf{Limitación relevante}\\
\hline
FAQ o árbol de diálogo & Comportamiento predecible para rutas predefinidas. &
Escasa flexibilidad ante paráfrasis y mantenimiento costoso de rutas.\\
Modelo generativo sin conocimiento controlado & Interacción flexible y lenguaje
natural. & Riesgo elevado de información no verificable o desactualizada.\\
Agente generativo con recuperación y gobierno & Puede unir consultas variadas con
fuentes seleccionadas y trazabilidad. & Depende de la calidad de la recuperación,
de las instrucciones y de controles adicionales.\\
\hline
\end{tabular}
\end{table}

La contribución propuesta del TFM es un artefacto de alcance restringido que
combina una base de conocimiento curada, reglas explícitas de seguridad,
integración Angular y una batería de evaluación orientada a evidencia. Su valor
no reside en afirmar que un agente sustituye a especialistas, sino en mostrar
cómo puede diseñarse un primer nivel de apoyo que mantenga la trazabilidad,
reconozca sus límites y pueda evolucionar mediante pruebas reproducibles.

.
\chapter{Estado del arte}
\label{chap:estado-arte}

Este capítulo revisa los conceptos y tecnologías que fundamentan el proyecto:
agentes conversacionales con recuperación de conocimiento, plataformas de
orquestación, integración web y gobierno de la IA generativa. El objetivo no es
atribuir al artefacto capacidades que no hayan sido verificadas, sino justificar
las decisiones de diseño y situar sus límites.

\section{De los asistentes basados en reglas a los agentes fundamentados}
\label{sec:evolucion-agentes}

Los asistentes conversacionales clásicos suelen organizarse como árboles de
diálogo o conjuntos de intenciones. Esta estrategia puede resultar predecible
para procesos muy acotados, pero exige anticipar numerosas formulaciones y rutas
de conversación. En una consulta de ciberseguridad, una misma intención puede
expresarse de forma muy distinta: por ejemplo, un usuario puede preguntar por un
correo sospechoso, un enlace que no reconoce o un posible robo de credenciales.
El problema se agrava cuando la respuesta depende de una política que cambia, de
un procedimiento interno o del contexto de un incidente.

Los modelos de lenguaje permiten generar una respuesta a partir de lenguaje
natural sin definir cada formulación de forma manual. Esta capacidad facilita el
acceso a documentación heterogénea, pero no convierte al modelo en una fuente de
verdad. Los modelos generativos pueden producir contenido erróneo presentado con
confianza, un fenómeno denominado por NIST \emph{confabulación}
\cite{nistGenAi2024}. Para un dominio como la ciberseguridad, la calidad no debe
medirse sólo por la fluidez de la conversación: importa que los hechos estén
fundamentados, que las recomendaciones sean prudentes y que la respuesta no
amplíe una solicitud peligrosa.

La generación aumentada por recuperación o RAG incorpora una fuente documental
externa al proceso de generación. Lewis et al. describen arquitecturas que
combinan memoria paramétrica y no paramétrica para tareas intensivas en
conocimiento, con el objetivo de recuperar pasajes relevantes y condicionar la
generación sobre ellos \cite{lewis2020rag}. En una solución empresarial, la idea
central es usar documentos y repositorios controlados como contexto de respuesta,
con la ventaja de poder actualizar el conocimiento sin reentrenar el modelo.

No obstante, RAG no garantiza por sí solo veracidad, actualidad o seguridad. La
recuperación puede seleccionar material irrelevante; el material puede estar
desactualizado, ser ambiguo o contener instrucciones maliciosas; y el modelo
puede resumirlo incorrectamente. Por ello, el presente trabajo trata el
conocimiento como un activo gobernado: se selecciona, se versiona, se limita por
alcance y se valida mediante preguntas cuya respuesta pueda contrastarse con una
fuente aprobada.

\section{Conocimiento de ciberseguridad para el caso de uso}
\label{sec:conocimiento-ciberseguridad}

El agente no necesita reproducir un marco completo de gestión de seguridad, pero
sí debe encajar en él. CSF 2.0 ofrece un vocabulario útil para relacionar las
consultas de personas usuarias con actividades de gobierno, protección,
detección, respuesta y recuperación \cite{nistCsf2024}. Por ejemplo, una duda
sobre autenticación multifactor se conecta principalmente con la protección; una
notificación de correo malicioso debe llevar hacia detección y respuesta; y una
pregunta sobre responsables o procedimiento remite a gobierno.

El contenido de respuesta ante incidentes requiere un tratamiento especialmente
cuidadoso. NIST recomienda incorporar la respuesta a incidentes en las
actividades de gestión del riesgo durante todo el ciclo de vida
\cite{nistIncident2025}. La guía de INCIBE sobre respuesta a incidentes subraya
la necesidad de responsables, procedimientos de notificación, registro de
evidencias y revisión del plan \cite{incibeCrisis}. De estos documentos se
extrae una decisión de diseño: el agente puede explicar pasos iniciales de
seguridad y el mecanismo de escalado, pero no debe simular una investigación,
ordenar acciones irreversibles ni sustituir a la cadena de decisión humana.

Los recursos públicos también ayudan a fijar un lenguaje accesible. INCIBE
describe el phishing como una suplantación destinada a obtener información,
causar un perjuicio económico o infectar el dispositivo \cite{incibePhishing}.
Esta definición sirve como punto de partida para preguntas frecuentes, pero la
respuesta del agente debe centrarse en la conducta segura: no abrir enlaces o
adjuntos sospechosos, usar el canal de reporte configurado y no aportar
credenciales ni datos personales en el chat.

La actualidad es otro requisito. El informe ENISA Threat Landscape 2025 analiza
incidentes y tendencias dentro de un periodo concreto y evidencia que la
prioridad de amenazas depende del contexto temporal y sectorial
\cite{enisaThreat2025}. Por tanto, no se deberá presentar una fuente de panorama
de amenazas como una descripción permanente ni inferir que sus estadísticas se
aplican a una organización ficticia. Su función en el TFM es justificar el
contexto y orientar la selección de casos, no alimentar afirmaciones específicas
que no estén respaldadas.

\section{Microsoft Copilot Studio como plataforma de agente}
\label{sec:copilot-studio}

Microsoft Copilot Studio es una plataforma para crear, configurar, probar y
publicar agentes con fuentes de conocimiento, temas, herramientas y canales. La
documentación de producto describe que el conocimiento puede añadirse a nivel de
agente o de tema y que admite, según el entorno y la licencia, sitios públicos,
documentos, SharePoint, Dataverse, Azure AI Search y otros orígenes
\cite{copilotKnowledge}. Esta diversidad es útil, pero aumenta la importancia de
aplicar el principio de mínimo conocimiento: se incorporará únicamente el
material necesario para las categorías de consulta aprobadas.

La orquestación generativa introduce una capa de planificación basada en modelos
de lenguaje que interpreta la intención y puede seleccionar conocimientos,
herramientas o temas para construir una respuesta \cite{copilotOrchestration}.
Frente a un árbol de diálogo totalmente manual, esa aproximación reduce parte de
la rigidez conversacional; a cambio, el comportamiento depende de las
instrucciones, descripciones y recursos disponibles. La plataforma recomienda
que las instrucciones se apoyen en recursos configurados, empleen nombres
exactos cuando se refieran a herramientas y se prueben tras cada cambio
\cite{copilotInstructions}.

Para este TFM, las instrucciones se utilizarán como una capa de orientación, no
como un control de seguridad único. Deben establecer el dominio permitido, el
formato de respuesta, la obligación de indicar límites y las condiciones de
derivación. Las decisiones críticas ---como bloquear la publicación no
autenticada, filtrar endpoints o restringir conectores--- requieren además
configuración de plataforma, control de acceso y revisión humana.

Las fuentes públicas merecen una precaución adicional. La documentación de
Microsoft indica que la recuperación sobre sitios públicos depende de contenido
indexado y posee límites de profundidad y de cobertura \cite{copilotPublicWeb}.
En consecuencia, una URL añadida como fuente no equivale a haber validado que
todas sus páginas relevantes serán recuperadas. El inventario de conocimiento
debe acompañarse de pruebas de recuperación para las preguntas que se pretenda
soportar.

\section{Evaluación de agentes conversacionales}
\label{sec:evaluacion-agentes}

La evaluación de un agente puede abarcar respuesta única, conversaciones de
varios turnos, recuperación de conocimiento, llamadas a herramientas, seguridad
y experiencia de usuario. Copilot Studio permite organizar casos en conjuntos de
prueba y comparar la salida con respuestas esperadas o métodos de calidad
\cite{copilotEvaluation}. Esta capacidad facilita la regresión: un mismo conjunto
puede ejecutarse después de modificar una fuente, una instrucción o una
integración.

Sin embargo, una evaluación automática debe interpretarse con cautela. Una
coincidencia textual puede penalizar una respuesta correcta formulada de otra
manera; una medida semántica puede aceptar una respuesta que omite una acción
esencial; y una puntuación de calidad no demuestra que la fuente recuperada sea
la adecuada. Por ello, el protocolo del Capítulo~\ref{chap:objetivos-metodologia}
combina criterios automáticos, cuando estén disponibles, y revisión humana con
una rúbrica. También reserva casos específicos para abstención segura,
derivación y resistencia frente a entradas adversariales.

La trazabilidad aporta un criterio adicional. Una respuesta no debe citar una
fuente inexistente, una URL no relacionada o una política ficticia como si fuera
real. Cuando el canal permita mostrar referencias, se comprobará que la
referencia corresponda al contenido utilizado. Cuando no las muestre, la
trazabilidad deberá aportarse mediante registros de prueba y un mapeo manual
entre respuesta esperada y fuente de conocimiento.

\section{Integración web y seguridad de canal}
\label{sec:integracion-web}

Una aplicación Angular permite controlar el contexto visual, los avisos al
usuario, los enlaces a procedimientos y la experiencia de consulta. Angular
recomienda mantener las dependencias actualizadas, evitar APIs identificadas como
riesgo de seguridad y reducir el riesgo de XSS mediante plantillas, sanitización,
Content Security Policy y Trusted Types cuando resulte aplicable
\cite{angularSecurity}. No obstante, la protección frente a XSS no resuelve por
sí sola los problemas propios de una aplicación de chat: gestión de sesión,
autorización, protección de tokens, controles de origen, registros y límites de
uso deben diseñarse de forma conjunta.

Copilot Studio puede publicar agentes en diversos canales. Microsoft recomienda
aplicar autenticación y controles de gobierno, especialmente cuando el agente se
utiliza dentro de una organización o por un grupo específico de usuarios
\cite{copilotPublish}. Para una aplicación web personalizada, el diseño deberá
declarar la opción de autenticación seleccionada y justificarla. Una integración
sin autenticación puede ser apropiada para una demostración sin conocimiento
interno, pero no constituye el diseño por defecto de un despliegue organizativo.

La documentación de seguridad del canal web indica que Copilot Studio utiliza el
canal Direct Line de Bot Framework para conectar la página o aplicación con el
agente. Si se habilita el acceso protegido, el cliente debe usar secretos o
tokens; en código que se ejecuta en el navegador, el secreto debe intercambiarse
por un token desde un servicio y no exponerse en el cliente
\cite{copilotWebSecurity}. Esta propiedad condiciona la arquitectura: un secreto
de Direct Line no se incluirá en variables públicas de Angular ni en un
repositorio.

Microsoft también documenta opciones de integración mediante SDK y Direct Line,
incluido un patrón basado en eventos para enviar contexto inicial de forma fiable
\cite{copilotDirectLine}. Dicho contexto deberá limitarse a datos no sensibles y
estar sujeto a validación. Si se habilita inicio de sesión único, Copilot Studio
requiere una configuración de autenticación con Microsoft Entra ID, registros de
aplicación diferenciados y un flujo de intercambio de tokens
\cite{copilotSso}. Esta configuración se tratará como una mejora opcional
condicionada a los requisitos y recursos del entorno, no como una funcionalidad
que pueda declararse implantada sin evidencia.

La viabilidad operativa también depende de licencias y capacidad. Microsoft
distingue licencias de usuario para crear y gestionar agentes de la licencia de
tenant, y recomienda estimar el consumo antes de un despliegue
\cite{copilotLicensing}. En este trabajo, los aspectos de licencia se documentan
como una restricción de entorno. No se infieren costes ni derechos de uso que no
puedan verificarse en el tenant concreto.

\section{Riesgos de seguridad específicos de la IA generativa}
\label{sec:riesgos-ia-generativa}

El modelo de amenazas debe considerar tanto vulnerabilidades tradicionales de la
aplicación como riesgos propios del agente. El AI RMF de NIST organiza la gestión
de riesgos de IA alrededor de las funciones de gobernar, mapear, medir y
gestionar \cite{nistAiRmf2023}. Su perfil para IA generativa señala riesgos
relacionados con la confabulación, privacidad, seguridad de la información,
integración de componentes y configuración humano--IA \cite{nistGenAi2024}.

OWASP incluye entre los riesgos de aplicaciones con LLM la inyección de prompts,
la divulgación de información sensible, los riesgos de cadena de suministro, el
envenenamiento de datos y la agencia excesiva \cite{owaspLlm2026}. Estos riesgos
se traducen en controles concretos para el proyecto:

\begin{itemize}
  \item no incorporar a la base de conocimiento texto no revisado ni permitir que
  una instrucción incluida en un documento redefina el comportamiento del agente;
  \item no almacenar secretos, tokens, credenciales o evidencias reales en
  fuentes de conocimiento ni en las conversaciones de prueba;
  \item limitar las acciones del agente a orientación informativa; no se
  habilitarán herramientas con capacidad de modificar sistemas;
  \item validar que las salidas presentadas por la interfaz no se interpreten como
  HTML confiable ni como instrucciones ejecutables;
  \item incorporar avisos visibles sobre el alcance, la posible incertidumbre y
  la necesidad de recurrir a canales humanos ante incidentes; y
  \item probar explícitamente intentos de revelar instrucciones, alterar el
  contexto o conseguir información de otros usuarios.
\end{itemize}

La defensa no debe descansar sobre una única instrucción de rechazo. La reducción
de superficie se consigue combinando alcance limitado, conocimiento aprobado,
autenticación adecuada, seguridad de canal, mínimo privilegio, revisión de
configuración y pruebas de robustez. Las limitaciones que sobrevivan a estos
controles deben aparecer en la memoria y en la interfaz, en lugar de ocultarse
detrás de un lenguaje excesivamente seguro.

\section{Privacidad, regulación y supervisión humana}
\label{sec:privacidad-regulacion}

Un agente conversacional puede recibir datos personales o información sensible
aunque no se haya diseñado para ello. El RGPD exige que el responsable aplique
medidas técnicas y organizativas apropiadas al riesgo y pueda demostrar el
cumplimiento \cite{gdpr}. Para el TFM, el control más directo es no utilizar
información real en pruebas y minimizar los datos que el usuario puede aportar.
En un uso futuro, habría que analizar la base jurídica, información al usuario,
retención, transferencias, subencargados, derechos de los interesados y gestión
de brechas.

El Reglamento (UE) 2024/1689 establece reglas armonizadas sobre sistemas de IA y
persigue promover una IA centrada en las personas y fiable, protegiendo a la vez
salud, seguridad y derechos fundamentales \cite{aiAct}. La clasificación y las
obligaciones concretas dependen del caso de uso, rol y fecha de aplicación; no se
deduce automáticamente que este prototipo cumpla o incumpla una categoría
regulatoria. Sus principios sí informan la solución: transparencia acerca de la
interacción con un agente, supervisión humana, documentación y tratamiento
proporcional del riesgo.

\section{Síntesis y contribución propuesta}
\label{sec:sintesis-estado-arte}

La Tabla~\ref{tab:comparacion-enfoques} resume el espacio de diseño. No pretende
medir el rendimiento de una implementación concreta; identifica qué propiedad
aporta cada enfoque y qué limitación debe compensarse.

\begin{table}[H]
\centering
\caption{Comparación conceptual de enfoques de asistencia conversacional.}
\label{tab:comparacion-enfoques}
\begin{tabular}{p{0.24\textwidth}p{0.31\textwidth}p{0.34\textwidth}}
\hline
\textbf{Enfoque} & \textbf{Ventaja principal} & \textbf{Limitación relevante}\\
\hline
FAQ o árbol de diálogo & Comportamiento predecible para rutas predefinidas. &
Escasa flexibilidad ante paráfrasis y mantenimiento costoso de rutas.\\
Modelo generativo sin conocimiento controlado & Interacción flexible y lenguaje
natural. & Riesgo elevado de información no verificable o desactualizada.\\
Agente generativo con recuperación y gobierno & Puede unir consultas variadas con
fuentes seleccionadas y trazabilidad. & Depende de la calidad de la recuperación,
de las instrucciones y de controles adicionales.\\
\hline
\end{tabular}
\end{table}

La contribución propuesta del TFM es un artefacto de alcance restringido que
combina una base de conocimiento curada, reglas explícitas de seguridad,
integración Angular y una batería de evaluación orientada a evidencia. Su valor
no reside en afirmar que un agente sustituye a especialistas, sino en mostrar
cómo puede diseñarse un primer nivel de apoyo que mantenga la trazabilidad,
reconozca sus límites y pueda evolucionar mediante pruebas reproducibles.

